\section{XF-7: Laguerre hierarchy and signed transverse tail}
\label{sec:xf7-laguerre-signed-tail}

XF-7 refines the transverse-mass programme by retaining oscillatory
cancellation in the $b$ coordinate.  The starting point is the standard
Riemann kernel $\Phi$ together with the transformed coordinate
\begin{equation}
s(t)=\Phi(\sqrt t),\qquad t>0.
\end{equation}
Planat and Solé study the first Laguerre quantity
\begin{equation}
F(t)=s'(t)^2-s(t)s''(t)
\end{equation}
and report strict log-concavity of $F$ together with the associated double
Turán hierarchy \cite{PlanatSole2026Second}.  In this work that result is
tracked as an external preprint claim.  The coordinate identities below are
exact consequences of the chain rule.

Put $t=r^2$. Then
\begin{equation}
s'(t)=\frac{\Phi'(r)}{2r},
\qquad
s''(t)=\frac{\Phi''(r)}{4r^2}-\frac{\Phi'(r)}{4r^3},
\end{equation}
so
\begin{equation}
F(r^2)=
\frac{r\bigl(\Phi'(r)^2-\Phi(r)\Phi''(r)\bigr)
+\Phi(r)\Phi'(r)}{4r^3}.
\label{eq:xf7-F-crosswalk}
\end{equation}
Define
\begin{equation}
A_0(r)=-\frac{\Phi'(r)}{\Phi(r)},
\qquad
B_0(r)=rA_0'(r)-A_0(r).
\end{equation}
Equation \eqref{eq:xf7-F-crosswalk} becomes
\begin{equation}
\boxed{
F(r^2)=\frac{\Phi(r)^2 B_0(r)}{4r^3}.
}
\label{eq:xf7-b0-crosswalk}
\end{equation}
The classical kernel is positive and strictly decreasing for $r>0$; this
property is used as standard external input \cite{Planat2026Theta}.  Hence
$A_0(r)>0$.  If $F(r^2)>0$ is supplied, then \eqref{eq:xf7-b0-crosswalk}
gives $B_0(r)>0$.

Now set
\begin{equation}
q(r)=\frac{A_0(r)}{r}.
\end{equation}
Direct differentiation yields
\begin{equation}
\boxed{
q'(r)=\frac{B_0(r)}{r^2}.
}
\label{eq:xf7-qprime}
\end{equation}
Therefore $B_0>0$ makes $q$ strictly increasing.  Moreover,
\begin{equation}
B_0=rA_0'-A_0>0
\quad\Longrightarrow\quad
A_0'>\frac{A_0}{r}>0,
\end{equation}
and consequently
\begin{equation}
(\log\Phi)''=-A_0'<0.
\end{equation}
Thus the first-Laguerre positivity in the $s(t)=\Phi(\sqrt t)$ coordinate,
combined with the standard radial decrease of $\Phi$, supplies an exact
conditional bridge to direct strict log-concavity of $\Phi$.

For the XF-6 transverse mass
\begin{equation}
M(a,b)=\Phi(a+b)\Phi(a-b),
\qquad 0\le b<a,
\end{equation}
one has
\begin{equation}
-\partial_b\log M(a,b)
=A_0(a+b)-A_0(a-b).
\end{equation}
Since $A_0(r)=rq(r)$ and $q$ is positive and increasing,
\begin{equation}
-\partial_b\log M(a,b)
\ge 2b\,q(a-b).
\end{equation}
Integration gives the adaptive envelope
\begin{equation}
\boxed{
M(a,b)
\le M(a,0)
\exp\!\left[-b^2\frac{A_0(a-b)}{a-b}\right].
}
\label{eq:xf7-adaptive-envelope}
\end{equation}
Its exponent is generated by the kernel itself rather than by an externally
selected curvature floor.

The main large-frequency gain comes from preserving the cosine.  If
$b\mapsto M(a,b)$ is nonincreasing on $[r,a]$ and $x\ne0$, one integration by
parts yields
\begin{equation}
\boxed{
\left|
\int_r^a M(a,b)\cos(2xb)\,\dd b
\right|
\le \frac{M(a,r)}{|x|}.
}
\label{eq:xf7-signed-tail}
\end{equation}
Combining \eqref{eq:xf7-adaptive-envelope} and
\eqref{eq:xf7-signed-tail} gives
\begin{equation}
\left|
\int_r^a M(a,b)\cos(2xb)\,\dd b
\right|
\le
\frac{M(a,0)}{|x|}
\exp\!\left[-r^2\frac{A_0(a-r)}{a-r}\right].
\label{eq:xf7-adaptive-signed-tail}
\end{equation}
This estimate retains an explicit $|x|^{-1}$ gain from phase cancellation.
It directly controls the transverse-oscillatory part of the XF-5 kernel.

The remaining correlated obligation is the sector containing
\begin{equation}
b^2\cos(2xa)\cosh(2yb),
\end{equation}
whose oscillation is longitudinal in $a$.  The global signed core--tail
balance is therefore recorded as an \emph{open sufficient route}.  The typed
XF-7 solver admits only the forward implication from that global gate to the
Wiener--Laguerre positivity criterion.  RH remains OPEN.
